\chapter{Methodology}\label{chap:methodology}

This chapter first introduces the concepts and metrics used (Section~\ref{sec:fundamentals}), then
describes the approach: data, retrievers, generators, judge and abstention rule
(Section~\ref{sec:approach}), and finally how it was implemented
(Section~\ref{sec:implementation}).

\section{Fundamentals}\label{sec:fundamentals}

\subsection{Retrieval-Augmented Generation}
A RAG system answers a question $q$ in two steps. A \emph{retriever} selects the $k$ passages from
a corpus that best match $q$, and a \emph{generator}, here a language model, receives $q$ and these
passages and produces an answer $a$. The answer is only as well founded as the passages allow, so
this thesis measures the two steps separately.

\subsection{Multi-hop Questions}
A \emph{multi-hop} question needs facts from more than one passage. In HotpotQA
\cite{yang2018hotpotqa}, every question used here has two \emph{gold} passages. A \emph{bridge}
question reaches the answer through an entity connecting the two passages, and a \emph{comparison}
question compares two entities. Table~\ref{tab:examples} shows one example of each from the data.
Multi-hop questions suit this study because retrieval can succeed partly: it can find one gold
passage and miss the other.

\begin{table}[htbp]
\centering\small
\caption{One question of each type, taken from the data used in this thesis.}\label{tab:examples}
\begin{tabular}{@{}>{\raggedright\arraybackslash}p{2.2cm}>{\raggedright\arraybackslash}p{6.5cm}%
                >{\raggedright\arraybackslash}p{1.7cm}>{\raggedright\arraybackslash}p{4cm}@{}}
\toprule
Type & Question & Answer & Gold passages \\
\midrule
Bridge & What was Iqbal F. Qadir on when he participated in an attack on a radar station located on
western shore of the Okhamandal Peninsula? & flotilla & Dwarka; Iqbal F. Qadir \\
Comparison & What is the shared country of ancestry between Art Laboe and Scout Tufankjian? &
Armenian & Art Laboe; Scout Tufankjian \\
\bottomrule
\end{tabular}
\end{table}

\subsection{Retrieval Metrics}
Let $G$ be the set of two gold passages of a question and $R_k$ the top $k$ retrieved passages.
Each metric below is averaged over the questions.
\begin{align}
\text{Recall@}k &= \frac{|G \cap R_k|}{|G|}, \label{eq:recall}\\
\text{Both@}k &= \mathbb{1}\left[\,G \subseteq R_k\,\right], \label{eq:both}\\
\text{MRR} &= \frac{1}{\text{rank of the first gold passage}}. \label{eq:mrr}
\end{align}
MRR is computed on the ten saved passages and is 0 if none of them is gold. Both@5 is the key value
for this study, because it is the share of questions for which the generator sees all the evidence
it needs.

\subsection{Answer Accuracy}
\textbf{Exact Match (EM)} and \textbf{F1} follow the official HotpotQA script
\cite{yang2018hotpotqa}. Answer and gold answer are normalised (lower case, no punctuation, no
articles). EM is 1 if they are identical. F1 is the harmonic mean of token precision and recall. If
either side is ``yes'', ``no'' or ``noanswer'' and the two differ, F1 is 0. An answer counts as
\emph{correct} if $\text{F1} \geq 0.5$ and the model did not abstain.

\subsection{Faithfulness and Hallucination}
Accuracy and faithfulness are different questions: an answer can be true but absent from the
passages, or supported by a passage and still wrong. Faithfulness asks only whether the passages
back the answer. Following RAGTruth \cite{niu2024ragtruth}, each answer gets one of three labels:
\begin{itemize}
\item \textbf{supported}: the context states or clearly implies the answer;
\item \textbf{contradicted}: the context says something that conflicts with the answer;
\item \textbf{unsupported}: the context does not contain enough information to back the answer.
\end{itemize}
The \textbf{hallucination rate} is the share of non-abstained answers that are labelled contradicted
or unsupported.

\subsection{Agreement}
Agreement between two raters is reported as the percentage of equal labels and as \textbf{Cohen's
kappa} \cite{cohen1960kappa}, which corrects for chance agreement:
\begin{equation}
\kappa = \frac{p_o - p_e}{1 - p_e},
\end{equation}
where $p_o$ is the observed agreement and $p_e$ the agreement expected by chance from the label
frequencies of the two raters. It is computed for the three labels, and for the binary split into
supported and hallucinated.

\subsection{Abstention and its Evaluation}
A system that may decline to answer has a trade-off between how often it answers and how good its
answers are \cite{geifman2017selective}. \textbf{Coverage} is the share of questions answered. For a
threshold $\tau$ on a confidence value, an answer is kept only if its confidence is at least $\tau$.
The \textbf{AUROC} of a confidence signal is the probability that a randomly chosen non-hallucinated
answer has higher confidence than a randomly chosen hallucinated one; 0.5 means that the signal is
useless and 1.0 that it separates the two perfectly.

\section{Approach}\label{sec:approach}

\subsection{Data}
The data come from the distractor configuration of HotpotQA \cite{yang2018hotpotqa}, validation
split, where all questions have the level ``hard''. Using seed \val{data.seed},
\val{data.questions} questions are sampled (\val{data.bridge} bridge, \val{data.comparison}
comparison). Retrieval is evaluated on all of them. Generation and judging use the first
\val{data.eval_questions} (\val{data.eval_bridge} bridge, \val{data.eval_comparison} comparison),
because of the limits of the free APIs (Section~\ref{sec:api-limits}).

The corpus holds the paragraphs of these \val{data.questions} questions and those of
\val{data.noise} further questions, which are never asked and act only as noise. This makes retrieval
harder and more realistic than searching only the ten paragraphs that come with each question.
Table~\ref{tab:data} gives the counts.

\begin{table}[htbp]
\centering
\caption{Data used in this thesis (counted from the processed files).}\label{tab:data}
\begin{tabular}{@{}lr@{}}
\toprule
Item & Count \\
\midrule
Questions with retrieval evaluation & \val{data.questions} \\
Questions with generation and judging & \val{data.eval_questions} \\
Additional questions used only as a source of noise passages & \val{data.noise} \\
Unique passages in the corpus & \val{data.passages} \\
Mean passage length (words) & \val{data.words} \\
Gold passages per question & \val{data.gold_per_q} \\
\bottomrule
\end{tabular}
\end{table}

\subsection{Retrievers and Conditions}
Four retrievers are compared. Each returns its top passages, and the generator receives the top
$k=\val{cfg.top_k}$ in that order.
\begin{description}[style=nextline,leftmargin=1em]
\item[BM25] Keyword search \cite{robertson2009bm25} over lower-cased word tokens.
\item[Dense] Semantic search \cite{karpukhin2020dpr} with the embedding model
  \texttt{BAAI/bge-small-en-v1.5} \cite{xiao2024cpack} and an exact FAISS inner-product index
  \cite{johnson2021faiss}. The embeddings are normalised, so the inner product equals the cosine
  similarity.
\item[Hybrid] Reciprocal rank fusion \cite{cormack2009rrf} of the top 100 passages from BM25 and
  from dense search, with constant $\val{cfg.rrf_k}$ (Fig.~\ref{fig:lst-hybrid}).
\item[Hybrid + rerank] The top \val{cfg.rerank} hybrid passages are scored again by the cross-encoder
  \nolinkurl{cross-encoder/ms-marco-MiniLM-L-6-v2} \cite{reimers2019sbert,nogueira2019rerank}, which
  reads the question and passage together. These scores are also the confidence signal of the
  abstention rule.
\end{description}
Two further conditions mark the limits. In the \textbf{closed-book} condition the model gets no
passages, which shows what it can answer alone. In the \textbf{oracle} condition it gets only the two
gold passages, which simulates perfect retrieval. Hallucination is not computed for the closed-book
condition, because there is no context to check the answer against.

\subsection{Generators}
Two small open-weight models are used through the Groq API: \texttt{openai/gpt-oss-20b}
\cite{openai2025gptoss} with \texttt{reasoning\_effort} low, and \texttt{qwen/qwen3.8-27b}
\cite{qwen2026qwen38} with \texttt{reasoning\_effort} none. Both use temperature
\val{cfg.temperature} and at most \val{cfg.max_tokens} tokens. The limit has to be this high because
both models can spend tokens on reasoning before they answer, and a low limit produced empty, cut-off
answers.

The prompts are in Fig.~\ref{fig:lst-prompt}. With passages, the model is told to answer using only
the context and to reply with a short answer. The prompt does not offer ``I don't know'' as an
option, so any refusal is the model's own. Such answers are detected with a list of phrases and count
as abstained.

\begin{figure}[htbp]
\lstinputlisting{generated/lst_prompt.py}
\caption{Prompts of the generator (\texttt{generation/llm.py}, \val{lst.prompt}).}\label{fig:lst-prompt}
\end{figure}

\subsection{Faithfulness Judge}\label{sec:judge}
Every non-abstained answer outside the closed-book condition is labelled by a language-model judge.
The judge is adapted from the faithfulness metric of RAGAS \cite{es2024ragas}. The answers are short,
so the whole answer is judged in one call instead of being split into statements. The judge sees the
same passages as the generator and is told to use only the context, not its own knowledge
(Fig.~\ref{fig:lst-judge}). No answer is marked as supported simply because its text appears in the
passages, since a model can pick the wrong name from a passage.

The main judge is \texttt{gemma-4-31b-it} (Google) \cite{gemma2026gemma4}, from a different model family
than both generators. The second judge, \texttt{openai/gpt-oss-120b} (Groq) \cite{openai2025gptoss}, is used
only to measure agreement, as it shares a family with one generator. All main results use one judge; the two
are never mixed in a table.

To check the judges, a balanced sample of \val{cfg.n_manual} answers was drawn in random order: 50 that the
main judge called hallucinated and 50 that it called supported. The author labelled all of them from
the passages only, in a tool that hid the judges' labels (Section~\ref{sec:res-judge}). The author was not fully
naive: before labelling, the author knew the pooled results and the sampling scheme, and the last 20 answers
were labelled after the results of the first 80 were known.

\begin{figure}[htbp]
\lstinputlisting{generated/lst_judge.py}
\caption{Prompt of the faithfulness judge (\texttt{eval/judge.py}, \val{lst.judge}).}\label{fig:lst-judge}
\end{figure}

\subsection{Abstention Rule}\label{sec:abstention}
The rule of RQ2 needs no model training and no change to the generator (Fig.~\ref{fig:abstention}). It is
tested on the hybrid + rerank and dense conditions with three confidence signals: \textbf{top1}, the retriever or
reranker score of the best passage; \textbf{top2}, the score of the second-best passage, since HotpotQA questions
need two passages and a weak second passage may indicate that a hop is missing; and \textbf{answer\_overlap}, the
share of the answer's normalised words found in the passages shown. ``Yes'' and ``no'' answers get 1.0 and are never
withheld, since those words say nothing about support. If the signal is below a threshold $\tau$, the system
withholds the answer.

\begin{figure}[htbp]
\centering
\begin{tikzpicture}[node distance=0.6cm and 0.7cm,
  box/.style={draw=ruleBrown,fill=boxFill,rounded corners=2pt,align=center,font=\small,
              minimum height=1.05cm,text width=2.6cm},
  arr/.style={-{Stealth},ruleBrown,thick}]
\node[box] (q) {Question};
\node[box,right=of q] (r) {Retrieval and generation of the answer};
\node[box,right=of r] (c) {Signal (top1, top2 or answer\_overlap) $\geq \tau$?};
\node[box,above right=0.1cm and 0.7cm of c,text width=2.4cm] (a) {The answer is given};
\node[box,below right=0.1cm and 0.7cm of c,text width=2.4cm] (s) {``I don't know''};
\draw[arr] (q) -- (r);
\draw[arr] (r) -- (c);
\draw[arr] (c.east) -- node[above,font=\scriptsize,sloped]{yes} (a.west);
\draw[arr] (c.east) -- node[below,font=\scriptsize,sloped]{no} (s.west);
\end{tikzpicture}
\caption{The abstention rule. In this study it is applied after generation to the saved answers, so
it needs no new API calls. With top1 or top2, a deployed system could skip the model call when it
abstains.}\label{fig:abstention}
\end{figure}

To avoid choosing the threshold on the answers that it is evaluated on, the answers are split
50/50 at random into a development half and a test half. A split this small is noisy, so the results
are the mean over 200 random splits. The threshold is chosen on the development half to maximise
\begin{equation}
\text{score} = \frac{\#\text{correct} - \#\text{hallucinated}}{n},
\end{equation}
where a withheld answer counts as neither (ties go to the higher coverage), and it is then reported
on the test half (Fig.~\ref{fig:lst-threshold}). Accuracy after abstention counts a withheld answer
as wrong. The AUROC of each signal is computed on all answers of a model, not only on the test half,
as it involves no fitting. The rule is training-free but not label-free: choosing $\tau$ needs about 50
labelled answers.

\begin{figure}[htbp]
\lstinputlisting{generated/lst_threshold.py}
\caption{Metrics at a threshold and the choice of the threshold (\texttt{run\_abstention.py},
\val{lst.threshold}). The repeated splits use the same logic in \texttt{run\_analysis.py}.}\label{fig:lst-threshold}
\end{figure}

\subsection{Further Analyses}
\begin{itemize}
\item \textbf{Retrieval against hallucination.} For each answer, the number of gold passages among the
  passages that the model saw (0, 1 or 2) is joined with its label. This is the key evidence for RQ1.
\item \textbf{Correctness against faithfulness.} A $2\times2$ count of correct or incorrect against
  supported or hallucinated. ``Correct but hallucinated'' indicates that an answer may come from the
  model's memory, rather than from the passages.
\item \textbf{Question type.} EM and hallucination rate for bridge and comparison questions.
\item \textbf{Error analysis.} \val{cfg.n_errors} hallucinated answers from the BM25 and dense conditions
  are assigned to one of five causes (wrong passage retrieved, one hop missing, passage misread, answer
  from memory or judge error) by the author. The author saw the judge's label and reason and the gold
  answer, so this step is not blind.
\item \textbf{Same-question checks.} Whether questions with a missed hop are simply harder is checked in two
  ways: through the oracle rate on the same questions, and by comparing a question with itself where the
  retrievers found different numbers of gold passages.
\item \textbf{Uncertainty.} 95\,\% confidence intervals use a percentile bootstrap
  \cite{efron1993bootstrap} with \val{cfg.bootstrap} resamples. Within one condition each answer is a
  different question, so answers are resampled. Pooled conditions contain the same question several times, so
  for them whole questions are resampled as well.
\end{itemize}

\section{Implementation}\label{sec:implementation}
Every step of the pipeline is a script (Fig.~\ref{fig:pipeline}), and all settings are kept in one file,
\texttt{config.py}.

\begin{figure}[htbp]
\centering
\begin{tikzpicture}[node distance=0.55cm and 0.45cm,
  box/.style={draw=ruleBrown,fill=boxFill,rounded corners=2pt,align=center,font=\footnotesize,
              minimum height=1.9cm,text width=3.9cm},
  arr/.style={-{Stealth},ruleBrown,thick}]
\node[box] (a) {\textbf{1. Data}\\ \texttt{prepare\_data.py}\\ questions, corpus};
\node[box,right=of a] (b) {\textbf{2. Retrieval}\\ \texttt{run\_retrieval.py}\\ four retrievers, metrics};
\node[box,right=of b] (c) {\textbf{3. Generation}\\ \texttt{run\_generation.py}\\ six conditions, two models};
\node[box,below=of c] (d) {\textbf{4. Judging}\\ \texttt{run\_eval.py}\\ two judges, verdicts saved};
\node[box,left=of d] (e) {\textbf{5. Abstention}\\ \texttt{run\_abstention.py}\\ no API calls};
\node[box,left=of e] (f) {\textbf{6. Analysis}\\ \texttt{run\_analysis.py}\\ no API calls};
\draw[arr] (a) -- (b);
\draw[arr] (b) -- (c);
\draw[arr] (c) -- (d);
\draw[arr] (d) -- (e);
\draw[arr] (e) -- (f);
\end{tikzpicture}
\caption{Overview of the experiment pipeline. The saved outputs of each step are the input of the
next, so abstention and analysis can be rerun without calling a model.}\label{fig:pipeline}
\end{figure}

\subsection{Software and Hardware}
Software: Python \val{env.python}, sentence-transformers \val{env.sentence_transformers},
\texttt{faiss-cpu} \val{env.faiss}, \texttt{rank\_bm25} \val{env.rank_bm25}, pandas \val{env.pandas},
scikit-learn \val{env.sklearn} and matplotlib \val{env.matplotlib}. The models are called through the
Groq and Google GenAI APIs. Hardware: a laptop with an 11th Gen Intel Core i7-11800H processor
(2.30\,GHz) and 16\,GB of RAM. It also has an NVIDIA GeForce RTX 3060 Laptop GPU (6\,GB), but the local
steps ran on the CPU, and the language models ran on the providers' servers.

\subsection{Retrieval}
BM25 and FAISS search run locally. The passage embeddings are computed once and stored, and the top ten
passages with their scores are saved per question, so retrieval does not have to be repeated. Fusion
of the two rankings is implemented in Fig.~\ref{fig:lst-hybrid}.

\begin{figure}[htbp]
\lstinputlisting{generated/lst_hybrid.py}
\caption{Hybrid retrieval by reciprocal rank fusion (\texttt{retrieval/retrievers.py},
\val{lst.hybrid}).}\label{fig:lst-hybrid}
\end{figure}

\subsection{Working within API Limits}\label{sec:api-limits}
The free tiers shaped how the experiment was run, and for that reason the number of questions:
\begin{itemize}
\item The Groq free tier allows 200,000 tokens per day and 8,000 per minute for each model. Each
  model has its own quota, so the models were run in parallel in separate processes.
\item The Gemma judge often returns server errors (HTTP 500 and 503), which are retried with a waiting
  time that doubles up to 60 seconds. Replies with no readable label are stored without a verdict and
  excluded (19 answers, Section~\ref{sec:res-main}).
\item Every answer and verdict is written to disk at once, one line per item. Rerunning a command skips
  completed work, so a run that stops at a limit continues where it stopped.
\item The \texttt{ragas} package did not work with the installed LangChain version, so the judge is
  implemented directly in \texttt{eval/judge.py}.
\end{itemize}

\subsection{Reproducibility}
The question sample, the development/test split and the bootstrap use fixed seeds. Generation uses
temperature \val{cfg.temperature}. All answers and verdicts are stored, so every later number can be
recomputed without calling a model. The tables in Chapter~\ref{chap:results} are written by a script from
the saved results (\texttt{thesis/make\_assets.py}), and the numbers in the text were checked against those
files. A hosted
model can change without notice, so the same request may not return the same answer later; the stored
outputs are therefore the reference.
